% Generated draft: requires mathematical and source review.
\documentclass[11pt]{article}
\usepackage[margin=25mm]{geometry}
\usepackage{fontspec}
\setmainfont{Latin Modern Roman}
\usepackage{amsmath,amssymb,mathtools,unicode-math}
\setmathfont{Latin Modern Math}
\usepackage{enumitem,needspace,xurl,hyperref}
\setlength{\parindent}{0pt}
\setlength{\parskip}{5pt}
\setlength{\emergencystretch}{3em}
\newcommand{\N}{\mathbb{N}}
\newcommand{\Z}{\mathbb{Z}}
\newcommand{\Q}{\mathbb{Q}}
\newcommand{\R}{\mathbb{R}}
\newcommand{\C}{\mathbb{C}}
\newcommand{\F}{\mathbb{F}}
\newcommand{\A}{\mathbb{A}}
\newcommand{\PP}{\mathbb{P}}
\newcommand{\E}{\mathbb{E}}
\newcommand{\eps}{\varepsilon}
\newcommand{\dd}{\,\mathrm{d}}
\begin{document}
\section*{UnsolvedMath ID: 800144}
\textbf{Source problem code:} catalogue entry
\subsection{Definitions and mathematical statement}
The supplied record does not define the sign-change counting function \(W^\psi(T)\), its parameter \(T\), or the constant \(\gamma_1\). Consequently, it does not determine a complete mathematical formulation. Under the usual convention, the Chebyshev function is
\[
\psi(x)=\sum_{\substack{p\text{ prime},\ k\geq 1\\p^k\leq x}}\log p\qquad(x\geq 1),
\]
and \(\gamma_1\) commonly denotes the smallest positive imaginary part of a nontrivial zero of the Riemann zeta function \(\zeta(s)\). These interpretations require confirmation from the source; they do not supply the missing definition of \(W^\psi\).

Reading the source's unindexed limit inferior as \(T\to\infty\) through its intended, unspecified parameter domain, write formally
\[
L=\liminf_{T\to\infty}\frac{W^\psi(T)}{\log T}.
\]
The source attributes to Morrill, Platt, and Trudgian, improving work of Kaczorowski, the bound
\[
L\geq\frac{\gamma_1}{\pi}+1.867\cdot10^{-30}.
\]
The requested problem is to improve this lower bound, with an improvement conditional on the Riemann hypothesis also permitted. Here the Riemann hypothesis means that every nontrivial zero of \(\zeta(s)\) has real part \(1/2\). Suggested targets, presented separately, are
\[
\text{(a)}\quad L\geq1.01\frac{\gamma_1}{\pi},\qquad
\text{(b)}\quad L\geq2\frac{\gamma_1}{\pi},\qquad
\text{(c)}\quad L=+\infty.
\]
The assumptions of the quoted existing bound are not specified in the supplied excerpt.

As heuristic motivation, the source expects large-scale positive and negative intervals of \(\psi(x)-x\) to alternate at approximately the rate suggested by a repeated-averaging argument, while each detected transition consists of many sign changes on a smaller scale, by analogy with random walks. It therefore suggests growth of \(W^\psi(T)\), and related unspecified counting functions, closer to \(\sqrt T\), up to logarithmic factors, than to \(\log T\). This is qualitative motivation, not a specified asymptotic conjecture.
\subsection{Short English statement}
Improve the supplied lower bound for the limit inferior of \(W^\psi(T)/\log T\), allowing the Riemann hypothesis as an additional assumption. The suggested targets range from a factor of \(1.01\) times \(\gamma_1/\pi\) to an infinite limit inferior. A complete statement requires the source's definitions of \(W^\psi\), \(T\), and \(\gamma_1\).
\subsection{Sources}
\begin{itemize}
\item Local UnsolvedMath registry, permanent ID \(800144\), exact source code \texttt{catalogue entry}, title ``Comparative Prime Number Theory --- Problem 1.'' The supplied record labels the problem open; this label is supplied metadata, not an independently verified assessment.
\item Hamieh, Kadiri, Martin, and Ng, \emph{Comparative Prime Number Theory Problem List} (2024), source item ``Problem 1 (\#13)'': \url{https://arxiv.org/abs/2407.03530}. The supplied provenance records access on 2026-10-16 and extraction by \texttt{arxiv\_latex}; no independent source visit or verification is asserted here.
\item The excerpt attributes the displayed lower bound to Morrill, Platt, and Trudgian, reference label [MPT], improving work of Kaczorowski, reference label [K2]. Full bibliographic entries and the hypotheses of those results are not supplied. The description of this bound as the best improvement known to the source speaker is a source-specific awareness statement.
\end{itemize}
\end{document}
